\documentclass[12pt]{article}
% load packages
\usepackage{amssymb,amsmath,amsfonts,geometry,graphicx,color,setspace,natbib,hyperref}
\usepackage{kpfonts}
%\usepackage[osf,sc]{mathpazo}
%\usepackage{lipsum}
%\usepackage{newpxtext,newpxmath}
\usepackage[tableposition=top, labelfont=bf]{caption}
\usepackage{booktabs}
\usepackage[flushleft]{threeparttable}
%\usepackage{subcaption}
\usepackage{siunitx}
%\usepackage[capposition=top]{floatrow}
\usepackage[table,xcdraw]{xcolor}
\usepackage{placeins}
\usepackage{lineno}


\usepackage{natbib}
\bibliographystyle{plainnatsw}

\usepackage{tabularx} % in the preamble
\usepackage{dsfont}
\usepackage{wrapfig}
\usepackage[nameinlink,capitalise]{cleveref}
\newcommand{\crefrangeconjunction}{--}

\usepackage{caption}
\DeclareCaptionFont{xipt}{\bfseries\fontfamily{qag}\selectfont}
\usepackage[font=xipt, labelfont=bf, textfont=normalfont]{caption}

%\renewcommand{\thesubfigure}{\Alph{subfigure}}
%\captionsetup[subfigure]{font={bf,large}, labelformat=simple, skip=1pt, singlelinecheck=false}

\usepackage{siunitx}
\sisetup{
  per-mode=symbol,
  group-separator={,},
  group-four-digits
}
\DeclareSIUnit{\ugpcm}{\micro\gram\per\cubic\meter}


\onehalfspacing
\geometry{left=1.2in,right=1.2in,top=1.0in,bottom=1.0in}
%\interfootnotelinepenalty=\@M

% sections options
\usepackage{titlesec}
\titleformat*{\section}{\fontfamily{qag}\selectfont\bfseries\LARGE}
\titleformat*{\subsection}{\fontfamily{qag}\selectfont\bfseries\large}
\titleformat*{\subsubsection}{\fontfamily{qag}\selectfont\bfseries\normalsize}
\usepackage{color, colortbl}
\definecolor{Gray}{gray}{0.9}
% color options
\definecolor{cerulean}{rgb}{0.0, 0.48, 0.65}
\definecolor{vermilion}{rgb}{0.89, 0.26, 0.2}
\hypersetup{
    colorlinks,
    linkcolor={cerulean},
    citecolor={cerulean},
    urlcolor={cerulean}
}
% fix color of brackets for citations
\bibpunct{\textcolor{cerulean}{(}}{\textcolor{cerulean}{)}}{,}{a}{}{;}

% load abstract package
\usepackage{abstract}
\renewcommand{\abstractnamefont}{\bfseries\fontfamily{qag}\selectfont}
\usepackage{epigraph}
\setlength\epigraphrule{0pt}
%\newfloatcommand{capbtabbox}{table}[][\FBwidth]

%change space in enumerates
\newenvironment{senumerate}{\enumerate\addtolength{\itemsep}{-10pt}}{\endenumerate} %espace entre les enumerate 

\usepackage{pdflscape} % landscape page
\usepackage{changepage} % change margins for specific pages
\usepackage{sidenotes} 

\usepackage[T1]{fontenc}
\usepackage{inconsolata}

\usepackage[symbol]{footmisc}


\renewcommand{\thefootnote}{\fnsymbol{footnote}}


\makeatletter
\AtBeginDocument{
\renewcommand\NAT@open{\color{cerulean}(}}
\makeatother

%----------------------------------------------------------------------

% DOCUMENT STARTS

%----------------------------------------------------------------------

\begin{document}

%----------------------------------------------------------------------

% TITLE PAGE

%----------------------------------------------------------------------

\begin{titlepage}
\vspace*{-2cm}
\begin{flushleft}
\Huge\fontfamily{qag}\selectfont\textbf{Why Acute Health Effects of Air Pollution Could Be Inflated \footnote[2]{\textbf{Comments and suggestions are highly welcome.} We are very grateful to Hélène Ollivier and Jeffrey Shrader for their fantastic guidance on this project. Many thanks to Geoffrey Barrows, Tarik Benmarhnia, Marie-Abèle Bind, Sylvain Chabé-Ferret, Marion Leroutier, Quentin Lippmann, Thomas Piketty, as well as seminars participants at the SusDev Colloquium at Columbia, IPWSD, Benmarhnia's Lab,  M\&A's Lab, and FAERE for their feedbacks.}}
\end{flushleft}
\renewcommand{\thefootnote}{\arabic{footnote}}
\vspace*{0.1cm} 
\begin{flushleft}
{\fontfamily{qag}\selectfontf\large
Vincent Bagilet\footnote{Columbia University - SIPA, New York, US. Email: \url{vincent.bagilet@columbia.edu}}\\[3pt] 
Léo Zabrocki\footnote{Paris School of Economics and
École des Hautes Etudes en Sciences Sociales, Paris, France. Email: \url{leo.zabrocki@psemail.eu}}
\\ [0.3cm] 
\normalsize{June 1st, 2022} \\
}

\end{flushleft}


\begin{abstract}
\small
\noindent Accurate and precise measurements of the short-term effects of air pollution on health play a key role in setting air quality standards. Yet, statistical power calculations are rarely---if ever---carried out. We first collect estimates and standard errors of all available articles found in the epidemiology and economics literatures. We find that nearly half of them may suffer from a low statistical power and could thereby produce statistically significant estimates that are actually inflated. We then run simulations based on real data to identify which parameters of research designs affect statistical power. Despite their large sample sizes, we show that studies exploiting rare exogenous shocks such as transport strikes or thermal inversions could have a very low statistical power, even if effect sizes are large. Our simulation results indicate that the observed discrepancy in the literature between instrumental variable estimates and non-causal ones could be partly explained by the inherent imprecision of the two-stage least-squares estimator. We also provide evidence that subgroup analysis on the elderly or children should be implemented with caution since the average number of events for an health outcome is a major driver of power. Based on these findings, we build a series of recommendations for researchers to evaluate the design of their study with respect to statistical power issues.

\bigskip
\end{abstract}
\setcounter{page}{0}
\thispagestyle{empty}
\end{titlepage}
\pagebreak \newpage

% double spacing for the rest of the document
\doublespacing
% numbering lines
%\linenumbers


%------------------------------------------------------------------------------

% SECTION - INTRODUCTION 

%------------------------------------------------------------------------------

\section{Introduction}
     
    % HOOK: THE IMPORTANCE OF MEASURING AHEAP
    From extreme events such as the London Fog of 1952 to the development of sophisticated time-series analyses, a vast scientific literature in epidemiology has established that air pollution induces adverse health effects on the very short-term \citep{schwartz1994people, le2002short, bell2004time, di_association_2017, liu_ambient_2019}. Increases in the concentration of several ambient air pollutants have been found to be associated with small relative increases in the daily mortality and emergency admissions for respiratory and cardiovascular causes \citep{samet2000national, shah2015short, orellano2020short}. All this evidence led to the implementation of public policies such as air quality alerts to mitigate the acute effects of air pollutants. Accurate estimates of the short-term health effects of air pollution are therefore crucial as they directly inform public health policies.
    
    % SECOND PARAGRAPH: TWO PARADOXES
    With this objective in mind, researchers in economics and epidemiology have addressed the issue of unmeasured confounding variables with causal inference methods in the last decade \citep{dominici_best_2017, bind_causal_2019}. Newly obtained results confirm the acute health effects of air pollution \citep{schwartz_estimating_2015, schwartz_national_2018, deryugina_mortality_2019}. Yet, causal estimates are often larger than what would have been predicted by the standard epidemiology literature. For instance, in Panel A of \Cref{fig:intro}, we see that instrumental variable estimates of 9 studies in the causal inference literature based on this method are always larger than naive ordinary least squares estimates. This could arguably be explained by the fact that instrumental variable strategies remove omitted variable bias and reduce attenuation bias coming from classical measurement error in air pollution exposure. Panel B of \Cref{fig:intro} however suggests an alternative and complementary explanation. For the 29 papers using causal inference methods found in this literature, we plot the standardized estimates against the inverse of their standard errors, which is a proxy for a study's precision. Large effect sizes are only found in imprecise studies and the more precise the study, the smaller the effect size. The negative relationship between effect sizes and studies' precision has also been observed in fields such as medicine, psychology and economics \citep{button2013power, camerer2018evaluating, schafer2019meaningfulness}.
    
        \begin{figure}[h!]
        \caption{Naive versus Causal Estimates and the Deflation of Effect Sizes as Precision Increases.}
            \label{fig:intro}
        \centering
        \includegraphics[width=\linewidth]{0.introduction/graph_paper.pdf}
        \caption*{\footnotesize \textit{Notes}: In Panel A, standardized estimates and their associated 95\% confidence intervals are displayed for the 9 articles of the causal inference literature based on instrumental variable strategies and for which estimates from naive regressions are available. Triangles represent instrumental variable estimates with dots are naive regression estimates. In panel B, standardized estimates of the 29 articles of the causal inference literature are plotted against the inverse of the standard errors, which can been considered as a measure of precision. Both axes are on a log10 scale.}
    \end{figure}
    
     The variation in studies' statistical power could explain the origin of this negative relationship but also help understand why causal estimates are larger than those found in the epidemiology literature. Simply put, studies with low precision result in larger effect sizes \citep{ioannidis2008most, gelman_beyond_2014}. Their statistical power is low and, to be statistically significant, their estimates need to be large enough, at least 2 standard errors away from 0 at the 5\% significance level. Since statistically significant results are more likely to be published, some estimates found in the literature may be inflated as they would come from a non-representative sample of the estimates, those large enough to be statistically significant \citep{brodeur2016star, brodeur_methods_2020}. The consequences of low statistical power are not specific to studies on short-term health effects of air pollution but may be particularly salient in this literature where the signal-to-noise ratio is often low \citep{peng2006model, peng2008statistical}.
    
    % OUR RESEARCH QUESTION AND WHAT WE DO: 1 PARAGRAPH
    In this paper, we undertake the first empirical investigation to determine if studies on the short-term health effects of air pollution could be under-powered and thereby produce inflated estimates. We start tackling this question by gathering a unique corpus of about 600 studies based on associations and 29 articles that rely on causal inference. For each of these papers, we run statistical power calculations to assess whether the design of the study would be robust enough to confidently detect an effect size smaller than the observed estimate \citep{gelman_beyond_2014, ioannidis_power_2017, lu_note_2019, timm_retrodesign_2019}. Using real data from the US National Morbidity, Mortality, and Air Pollution Study \citep{samet2000national}, we then implement simulations to identify the characteristics of research designs that drive their statistical power and the inflation of statistically significant estimates \citep{black2019simulated, gelman2020regression, altoe_enhancing_2020}. 
            
    % OUR RESULTS: 2 PARAGRAPHS
    The results of our statistical power calculations show that research designs based on associations and causal inference methods are similarly prone to statistical power issues. Half of the studies in the two strands of the literature have a statistical power below 80\% to detect effect sizes that are only 25\% smaller than their observed estimates. In the standard epidemiology literature, under-powered studies could produce statistically significant estimates 2 times larger than true effect sizes. Our retrospective power calculations also highlight a wide heterogeneity in the robustness of articles with respect to statistical power issues. For example, if the true effect sizes are equal to the ones predicted by the standard epidemiology literature, the statistical power of studies using instrumental variable designs would range from 5\% to 64\%. In some studies, statistically significant estimates would be just 1.3 times larger than the true effect sizes, while in others, the inflation factor could be as high as 41.
    
    Our simulation results help understand why some research designs face statistical power issues. We first show that a very large number of observations is needed for all causal inference methods to reach a sufficient statistical power. Regression discontinuity designs based on air quality alerts rely on sample sizes that may be too small for statistically significant estimates not to be inflated. Second, we show that the use of public transport strikes or thermal inversions as exogenous shocks on air pollution could be problematic. These studies are based on rare events, which in some cases represent less than 1\% of the observations. The resulting statistical power is very low, around 15\%, and statistically significant estimates can exaggerate even large true effect sizes by a factor of 2.7. Third, we find that the average daily count of cases of an health outcome is a key driver of statistical power for all empirical strategies. Statistically significant estimates of the effects of air pollution on the elderly or children can be very inflated since health outcomes for these groups often have few daily cases.

    % CONTRIBUTIONS: 2 PARAGRAPHS
    Our article makes two contributions to the literature on the acute health effects of air pollution. First, as highlighted by the replication crises in medicine, psychology and experimental economics \citep{button2013power, open2015estimating, camerer2018evaluating}, there is a crucial need to evaluate the deficiencies of current statistical practices grounded in the null hypothesis significance testing framework \citep{ziliak2008cult, simonsohn2014p, smaldino2016natural, greenland2017invited, christensen2019transparent, amrhein2019statistical}. Our paper participates in the growing literature that uses retrospective power calculations to evaluate the plausibility of published findings \citep{ioannidis_why_2008, gelman_beyond_2014, smaldino2016natural, ioannidis_power_2017, ferraro_featureis_2020, stommes2021reliability}. To the best of our knowledge, this paper is the first to show how to carry out retrospective statistical power calculations for studies on air pollution and human health. We also provide the first evidence that under-powered studies are a real issue in this field.
    
    Second, except for standard models used in the epidemiology literature \citep{winquist2012power}, few statistical power analyses exist to help researchers improve their study designs \citep{bhaskaran2013time}. This paper is the first to give concrete recommendations to avoid statistical power issues for several research designs estimating the acute health effects of air pollution. Statisticians have long advocated the use of fake-data simulations to flexibly evaluate the inference properties of statistical models \citep{gelman_beyond_2014, vasishth2019embrace, altoe_enhancing_2020, gelman2020regression}. In our paper, we follow this advice but rely instead on real data since it is very complex to correctly simulate the relationships between ambient air pollution, weather parameters, calendar indicators and health outcomes. Our article is more closely connected to three recent articles evaluating the type I error rate and the lack of statistical power of several panel data models used to estimate the impacts of public policies on mortality outcomes \citep{schell2018evaluating, black2019simulated, griffin2020variation}. These simulations focus on event study designs and treatment effects happening on a medium to long time scale. On the contrary, our simulations gauge the capacity of reduced-form, instrumental variable and regression discontinuity designs to estimate very short-run effects in the context of high-frequency data.
    
    Finally, we strive to make our analyses fully and easily reproducible to help researchers implement retrospective power calculations and power simulations in their own studies. We use state-of-the-art literate programming to explain and render all coding procedures in nicely formatted HTML documents \citep{allaire2018distill}. All replication and supplementary materials are available on this \href{https://vincentbagilet.github.io/inference_pollution/}{website}.
     
    % ROADMAP
    In the following section, we implement a simple simulation exercise to show why statistically significant estimates exaggerate true effect sizes when studies have a low statistical power. In section \hyperlink{section.3}{3}, we present our retrospective analysis of the literature. In section \hyperlink{section.4}{4}, we detail our simulation procedure to replicate empirical strategies. We display the simulation results in section \hyperlink{section.5}{5} and provide specific guidance on study design in section \hyperlink{section.6}{6}.
    

%------------------------------------------------------------------------------

% SECTION - BACKGROUND ON STATISTICAL POWER, TYPE M AND S ERRORS 

%------------------------------------------------------------------------------

\section{Background on Statistical Power, Type M and S errors}

    In a seminal paper, \cite{gelman_beyond_2014} point out that researchers working in the null hypothesis significance testing framework are often unaware that "statistically significant" estimates suffer from a winner's curse in under-powered studies: these estimates can largely overestimate true effect sizes and can even be of the opposite sign. In this section, we implement a simple simulation exercise to illustrate these two counter-intuitive issues and explain why they could matter in studies on the acute health effects of ambient air pollutants.

% A FICTIONAL EXAMPLE
    \subsection{Illustrative Example}

    Imagine that a mad scientist is able to implement a randomized experiment to measure the short-term effects of air pollution on daily non-accidental mortality. The experiment takes place in a major city over the 366 days of a leap year. The scientist is able to increase concentration of particulate matter with a diameter below 2.5 $\si{\um}$ (PM$_{2.5}$) by 10 $\si{\ugpcm}$---a large shock equivalent to one standard deviation increase in the concentration of PM$_{2.5}$. Concretely, the scientist implements a complete experiment where they randomly allocate half of the days to the treatment group and the other half to the control group. They then measure the treatment effect of the intervention by computing the average difference in means between treated and control outcomes: the estimate for the treatment effect is equal to 4 additional deaths and is "statistically significant" at the 5\% level, with a \textit{p}-value of 0.04. The statistical significance of the estimate fulfills the scientist expectations, who immediately starts writing their paper. Had they not obtained a statistically significant estimate, they might not have submitted their result. 
    
    Unfortunately for the scientist, we know what the true effect of the experiment is since we created the data. In \Cref{tab:science_table}, we display the Science table where we observe the pair of potential outcomes of each day, $Y_{i}(W_{i}=0)$ and $Y_{i}(W_{i}=1)$ \citep{rubin_estimating_1974}. $Y_{i}$ represents a daily count of non-accidental deaths and $W_{i}$ the treatment assignment, which is equal to 1 for treated units and 0 otherwise. We first simulated the daily non-accidental mortality counts in the absence of treatment (i.e., the $Y(0)$ column of \Cref{tab:science_table}), by drawing 366 observations from a negative binomial distribution with a mean of 106 and a variance of 402. We chose the parameters to approximate the distribution of non-accidental mortality counts in a large European city. We then defined the counterfactual distribution of mortality by adding, on average, 1 extra death (i.e., the $Y(1)$ column of \Cref{tab:science_table}). 
    
            \input{1.background/science_table}
            
    \renewcommand*{\thefootnote}{\arabic{footnote}}
    \setcounter{footnote}{0}

    \noindent This treatment effect size represents approximately a 1\% increase in the mean of the outcome\footnote{Note that the magnitude of this hypothetical effect is higher than what has been found in a recent and large-scale study based on 625 cities. \cite{liu_ambient_2019} found that a 10 $\si{\ugpcm}$ increase in $\text{PM}_{2.5}$ concentration was associated with a 0.68\% (95\% CI, 0.59 to 0.77) relative increase in daily all-causes mortality.}. Following the fundamental problem of causal inference, the daily count of deaths the scientist observes is given by the equation: $Y_{i}^{\textnormal{obs}} = W_{i} \times Y_{i}(1) + (1-W_{i})\times Y_{i}(0)$. Treated units express their $Y_{i}(1)$ values and control units their $Y_{i}(0)$ values. 
    
    With a random assignment of the treatment, how is it possible that the statistically significant estimate found by the scientist can be 4 times larger than the true treatment effect size? Replicating many times the experiment can help understand why.
        
% DEFINING STATISTICAL POWER, TYPE M AND S ERRORS
    \subsection{Defining Statistical Power, Type M and S errors}

        \begin{figure}[h!]
            \caption{Estimates of the 10,000 Simulations.}
            \label{fig:estimates_experiments}
            \centering
            \includegraphics[width=\linewidth]{1.background/graph_experiment_results.pdf}
            \caption*{\footnotesize \textit{Notes}: In Panel A, blue and red dots represent the point estimates of the 10,000 iterations of the randomized experiment ran by the mad scientist. Red dots are statistically significant at the 5\% level while blue dots are not. The black solid line represents the true average effect of 1 additional death.}
        \end{figure}

    In \Cref{fig:estimates_experiments}, we plot the estimates of 10,000 iterations of the experiment. If there is a large variation in the effect size of estimates, the average is reassuringly equal to the true treatment effect of 1 additional death. We can however see that estimates close to the true effect size would not be statistically significant at the 5\% level. In a world without publication bias, we could be confident that several replications of this experiment would recover the true treatment effect. Unfortunately, researchers are---despite recent changes in scientific practices and editorial policies---not incited enough to publish replication exercises and not statistically significant estimates. In a world with publication bias, only statistically significant estimates would be made public. Out of the 10,000 estimates, about 800 are statistically significant at the 5\% level. The \textit{statistical power} of the experiment, which can be defined as the probability to reject the null hypothesis when there is actually an effect, is therefore equal to 8\%. The scientist was therefore very lucky to get a statistically significant estimate. 
    
    But with such a low statistical power, statistically significant estimates cannot be trusted anymore. Two metrics, the average type M (for magnitude) error and the probability to make a type S (for sign) error are useful to assess the negative consequences of lacking statistical power. First, we can evaluate by how much statistically significant estimates are inflated compared to the true treatment effect size by computing the average ratio of the absolute values of the statistically significant estimates over the true effect size \citep{gelman_beyond_2014}. With a statistical power of 8\%, the scientist would on average make a type M error equal to 5! Second, we can notice that a non-negligible fraction of statistically significant estimates are of the wrong sign in \Cref{fig:estimates_experiments}: this proportion is the probability of making a type S error  \citep{gelman_type_2000}. For this experiment, a statistically significant estimate has a 8\% probability of being of the wrong sign! 
    
    Thus, if the scientist would like to estimate the effect of the experiment through the prism of the statistical significance, they would need a larger number of observations:  statistical power would then rise and conversely type M and S error would shrink.
    
\subsection{Relevance for Studies on Acute Health Effects of Air Pollution}

    Type M and S errors are two concepts that highlight the danger of having too much confidence in statistically significant estimates when studies are under-powered. This issue is virtually absent from the literature but studies on the acute health effect of air pollution could be under-powered for several reasons. First, researchers work with observational data and often cannot control the sample size of their studies due to data availability. Very few guidance on the drivers of studies' statistical power actually exists \citep{winquist2012power, bhaskaran2013time}. Moreover, reaching a large statistical power could be challenging since estimated effect sizes are remarkably small and the modeling of high-frequency variations in daily mortality or emergency admission is difficult \citep{peng2006model, peng2008statistical}. Finally, we observe both in the standard epidemiology and the causal inference literatures a negative relationship between estimated effect sizes and studies' precision. It is important to investigate if this pattern could be explained by imprecise studies making type M errors \citep{ioannidis2008most, gelman_beyond_2014, ioannidis_power_2017, ferraro_featureis_2020}.
    
%------------------------------------------------------------------------------

% SECTION - RETROSPECTIVE DESIGN ANALYSIS OF THE LITERATURE

%------------------------------------------------------------------------------

\section{Retrospective Analysis of the Literature}
    
    In this section, we first explain how to implement a retrospective analysis of a study. Using different scenarios about the true effect sizes of studies found in the standard epidemiology and causal inference literatures, we then assess to what extent they could suffer from low statistical power issues.
    
% COMPUTING STATISTICAL POWER, TYPE M AND S ERRORS
    \subsection{How to Run a Retrospective Analysis}
    
    Based on the notations of \cite{zwet_significance_2021}, the statistical framework for a retrospective power analysis can be formalized as follows. We take the general case where we want to estimate, for a given research design, the causal effect $\beta$ of an air pollutant on an health outcome. We assume that we have a normally unbiased estimate $b$ of $\beta$ with a standard error $s$. If we knew the true value of $\beta$, we could first compute the statistical power of our study for rejecting the null hypothesis $H_{0}: \beta = 0$. It is defined as $\Phi(-1.96-\frac{\beta}{s})+1-\Phi(1.96-\frac{\beta}{s})$, where $\Phi$ is the cumulative function of the standard normal distribution. Given the power of our study, we could then compute the average magnitude of a type M error, that is to say the average inflation of statistically significant estimates at the 5\% level for the study. It can be expressed as $\mathbb{E}(\frac{|b|}{|\beta|} | \beta, s, |b|/s>1.96)$. Finally, we could compute the probability to make a type S error, that is to say the probability that the estimate is of the opposite sign of $\beta$. It is given by $\frac{\Phi(-1.96-\frac{\beta}{s})}{1-\Phi(1.96-\frac{\beta}{s}) + \Phi(-1.96 - \frac{\beta}{s})}$. Unfortunately, we never observe the true value of the causal estimand of interest. We therefore need to make guesses about the value of $\beta$ to compute the three previous metrics. \cite{gelman_beyond_2014} propose to run simulations to calculate them. In our project, we rely on the closed-form expressions derived by \cite{lu_note_2019} and their implementation in the \textbf{\textsf{R}} package \texttt{retrodesign} developed by \cite{timm_retrodesign_2019}.
    
    It is very important to keep in mind that the usefulness of a retrospective power analysis relies entirely on a credible guess of the true effect size a study is trying to estimate. As the true effect is never observed, researchers can have very different priors on its magnitude. They could therefore assess differently the extent to which a study risks to suffer from statistical power issues. To illustrate this tension, we provide below a case study showing how a scientific discussion about effect sizes arises with a retrospective analysis. 
    
    In a flagship publication, \cite{deryugina_mortality_2019} instrument PM$_{2.5}$ concentrations with wind directions to estimate its effect on mortality, health care use, and medical costs among the US elderly. They gathered 1,980,549 daily observations at the county-level over the 1999–2013 period; it is one of the biggest sample sizes in the literature. When the authors instrument PM$_{2.5}$ with wind direction, they find that “a 1 $\si{\ugpcm}$ (about 10 percent of the mean) increase in PM$_{2.5}$ exposure for one day causes 0.69 additional deaths per million elderly individuals over the three-day window that spans the day of the increase and the following two days”. The estimate's standard error is equal to 0.061. In \Cref{fig:deryugina}, we plot the statistical power, the inflation factor of statistically significant estimates and the probability that they are of the wrong sign as a function of hypothetical true effect sizes.
    
        \begin{figure}[h!]
            \caption{Power, Type M and S Errors Curves for \cite{deryugina_mortality_2019}.}
            \label{fig:deryugina}
        \centering\includegraphics[width=\linewidth]{2.literature_review/iv_retrodesign_example.pdf}
            \caption*{\footnotesize \textit{Notes}: In each panel, a metric, such as the statistical power, the exaggeration ratio or the probability to make a type S error, is plotted against the range of hypothetical effect sizes. The "IV" label represents the value of the corresponding metric for an effect size equal to \cite{deryugina_mortality_2019}'s two-stage least square estimate. The "Epidemiology" label stands for the estimate found in \cite{di_association_2017}, which is the epidemiology article most similar to \cite{deryugina_mortality_2019}. The "OLS" label corresponds to the estimate found by \cite{deryugina_mortality_2019} when the air pollutant is not instrumented.}
        \end{figure}
    
    The estimate found by \cite{deryugina_mortality_2019} represents a relative increase of 0.18\% in mortality. We labeled it as "IV" in \Cref{fig:deryugina}. Is this estimated effect size large compared to those reported in the standard epidemiology literature? We found a similar article to draw a comparison. Using a case-crossover design and conditional logistic regression, \cite{di_association_2017} find that a 1 $\si{\ugpcm}$ increase in PM$_{2.5}$ is associated with a 0.105\% relative increase in all-cause mortality in the Medicare population from 2000 to 2012. The effect size found by \cite{deryugina_mortality_2019} is larger than this estimate labeled as "Epidemiology" in \Cref{fig:deryugina}. If the estimate found by \cite{di_association_2017} was actually the true effect size of PM$_{2.5}$ on elderly mortality, the study of \cite{deryugina_mortality_2019} would have enough statistical power to perfectly avoid type M and S errors. Now, suppose that the true effect of the increase in PM$_{2.5}$ was 0.095 additional deaths per million elderly individuals---the estimate the authors found with a "naive" multivariate regression model. The statistical power would be 34\%, the probability to make a type S error could be null but the overestimation factor would be on average equal to 1.7. Even with a sample size of nearly 2 million observations, \cite{deryugina_mortality_2019} could make a non-negligible type M error if the true effect size was the naive ordinary least square estimate. Yet, the authors could argue that their instrumental variable strategy leads to a higher effect size as it overcomes unmeasured counfounding bias and measurement error. Besides, for effect sizes down to 0.182 additional deaths per million elderly individuals (a 0.05\% relative increase), their study has a very high statistical power and would not run into substantial type M error. A retrospective analysis is thus a very convenient way to think about the statistical power of a study to accurately detect alternative effect sizes.
    
    % EPIDEMIOLOGY LITERATURE
    \subsection{Standard Epidemiology Literature}
    
        % introductory paragraph
        Hundreds of papers have been published on the short-term health effects of air pollution in epidemiology, medicine and public health journals. A large fraction of articles are based on Poisson generalized additive models, which allow to flexibly adjust for the temporal trend of health outcomes and for non-linear effects of weather parameters. This literature spans over 20 years and has replicated analyses in a large number of settings, providing crucial insights on the acute health effect of air pollution. Advocates of causal methods would surely argue that these articles could suffer from omitted variable biases. Even if they may be more biased, we find it valuable to assess their potential statistical power issues and compare them with causal inference papers.

       To gather a corpus of relevant articles, we use the following search query on \href{https://www.ncbi.nlm.nih.gov/}{PubMed} and \href{https://www.scopus.com/home.uri}{Scopus} to select studies on the short-term health effects of air pollution:
               
        \begin{quote}
       \footnotesize\texttt{'TITLE(("air pollution" OR "air quality" OR "particulate matter" OR "ozone"', 'OR "nitrogen dioxide" OR "sulfur dioxide" OR "PM10" OR "PM2.5" OR', ' "carbon dioxide" OR "carbon monoxide")',  'AND ("emergency" OR "mortality" OR "stroke" OR "cerebrovascular" OR', '"cardiovascular" OR "death" OR "hospitalization")', 'AND NOT ("long term" OR "long-term")) AND "short term"'}
        \end{quote}

       \noindent We retrieve the abstracts of 1834 articles. We then extract estimates and confidence intervals from these abstracts using regular expressions (regex). We illustrate this procedure using one sentence of a randomly selected article from this literature review \citep{vichit-vadakan_public_2008}:
        
        \begin{quote}
         \footnotesize    “\texttt{The excess risk for non-accidental mortality was \textbf{\color{cerulean}1.3\% [95\% confidence interval (CI), 0.8–1.7]} per 10 $\si{\ugpcm}$ of PM10, with higher excess risks for cardiovascular and above age 65 mortality of \textbf{\color{cerulean}1.9\% (95\% CI, 0.8–3.0)} and \textbf{\color{cerulean}1.5\% (95\% CI, 0.9–2.1)}, respectively.}”
        \end{quote}
        
        \noindent Our algorithm detects phrases such as “95\% confidence interval (CI)” or “95\% CI” and looks for numbers directly before this phrase or after and in a confidence interval-like format. Using this method, we retrieve 2666 estimates from 784 abstracts. We then read these abstracts and filter out articles whose topic falls outside of the scope of our literature review. Our corpus is thus composed of 668 articles for which we detect 2155 estimates. Importantly, the set of articles considered is limited to those displaying confidence intervals and point estimates in their abstracts.

        Based on this subset of articles, we implement a retrospective analysis in which we check the overall sensitivity of studies for true effect sizes expressed as fraction of observed estimates. Without carefully reading each article, we cannot make more informed guesses about true effect sizes since estimates are expressed for different increases in air pollution concentration. We think that our rough approach is still valuable since a well-designed study should be able to detect effect sizes smaller than the estimated one. For instance, if we find that a study has a statistical power of 30\% when we assume that the true effect is 25\% lower than the measured estimate, it is likely that the study is not very robust to statistical power issues.
        
        Our results for the standard epidemiology literature are at first sight reassuring. If the true effect sizes of the studies were equal to 75\% of estimated coefficients, the median statistical power would be equal to 85\% and the median exaggeration factor would be only 1.1. At least 50\% of this literature does not seem to suffer from substantial statistical power issues since their power would be above 80\%. Type S error is not an issue for most articles. Yet, even if the measured effect was close to the true effect, a non-negligible proportion of articles would display low statistical power and presents a substantial risk of making a type M error. About 47\% of estimates would not reach the conventional 80\% statistical power threshold if the true effect was 75\% the size of the measured effect. For these under-powered studies, the average type M is 1.9 and the median 1.5. We also observe that the proportion of under-powered studies has been stagnating since the 1990s, revealing that practices regarding statistical power have not evolved over time.
        
        Finally, skeptic researchers could rightly complain that assuming for each study a true effect size equal to 75\% of the estimate is arbitrary. To overcome this criticism, we expand our review of the standard epidemiology literature by running statistical power calculations based on two recent meta-analyses: one by \cite{shah2015short} on mortality and emergency admission for stroke, and the other one by \cite{orellano2020short} on broader causes of mortality. We use the meta-analysis estimates as true effect sizes for the 290 studies gathered by \cite{shah2015short} and \citep{orellano2020short}. This is the approach recommended by \cite{gelman_beyond_2014} and \cite{ioannidis_power_2017} to make more informed guesses about true effect sizes. 60\% of studies in \cite{orellano2020short} have a statistical power below 80\%. The median exaggeration ratio of statistically significant estimate is equal to 2. The proportion of under-powered studies is similar in \cite{shah2015short} but the median type M error is equal to 3. With more informed guesses about true effect sizes, we clearly see that under-powered studies are an issue in the standard epidemiology literature.
        
% CAUSAL INFERENCE LITERATURE
    \subsection{Causal Inference Literature}
    
        Using \href{https://scholar.google.com/}{Google Scholar} and \href{https://pubmed.ncbi.nlm.nih.gov/}{PubMed}, we search papers using causal inference methods and investigating the short-term effects of air pollution on mortality or emergency admission outcomes. Specifically, we only consider articles that exploit short-run exogenous shocks such as air pollution alerts, public transport strikes, changes in wind direction, thermal inversions, to name but a few. For instance, we did not select articles on the impact of low emission or congestion pricing zones as they evaluate health effects over several months or years. In \Cref{tab:lit_review}, we display the 29 articles that match our search criteria. We read each article and retrieve the estimates and standard errors for the main results: for simplicity, we only select one of the main results discussed by the researchers. We also record the numbers of observations and summary statistics on the outcome and independent variables to compare studies by standardizing the estimated effect sizes. 
    
%-----------------------------------------------------------------------
% TABLE FOR THE LITERATURE REVIEW   
% ATTENTION, C'EST UN PEU DE LA TAMBOUILLE POUR LA METTRE EN FORME
    \newgeometry{margin=1cm}
    \begin{landscape}
    \begin{table}[]
\thispagestyle{empty}
\caption{Our Corpus of Papers from the Causal Inference Literature.}
\begin{adjustwidth}{-0.5cm}{0cm}
\centering
\scriptsize
\begin{threeparttable}
\label{tab:lit_review}
\begin{tabular}{p{2.5cm}p{3cm}p{7cm}p{7cm}p{4cm}}
\toprule
\textbf{Article}      & \textbf{Location}           & \textbf{Health Outcome}                                     & \textbf{Independent Variables}                                                     & \textbf{Study Design}                             \\ \midrule
\rowcolor{Gray}
\cite{arceo_does_2016}    & Mexico City, Mexico         & Infant Mortality                                            & PM10, Thermal Inversion (IV)                                       & Instrumental Variable                             \\
\cite{austin2020covid}          & Counties, USA               & Rates of Confirmed COVID-19 Deaths                          & PM2.5 (air pollutant), Wind Direction (IV)                                         & Instrumental Variable                             \\
\rowcolor{Gray}
\cite{baccini_assessing_2017}         & Milan, Italy                & Non-Accidental Mortality                                    & Dummy for PM10 Concentration \textgreater To 40 µg/m³                              & Propensity Score Matching                         \\
\cite{barwick_morbidity_2018}         & All Cities, China           & Number of Health Spending Transactions                      & PM2.5, Spatial Spillovers of PM2.5 (IV)                            & Instrumental Variable                             \\
\rowcolor{Gray}
\cite{bauernschuster_when_2017}  & 5 Largest Cities, Germany   & Admissions for Abnormalities of Breathing (age below 5)     & PM10, Public Transport Strikes Dummy                               & Difference in Differences                         \\
\cite{beard_winter_2012}          & Salt Lake County, USA       & Emergency Visits For Asthma                                 & Thermal Inversions                                                                 & Time-stratified case-crossover design             \\
\rowcolor{Gray}
\cite{chen_effect_2018}           & Toronto, Canada             & Asthma-Related Emergency Department Visits                  & Air Quality Eligibility, Air Quality Altert                                        & Fuzzy Regression Discontinuity                    \\
\cite{deryugina_mortality_2019}      & Counties, USA               & All Causes of Mortality (Age 65+)                           & PM2.5, Wind Direction (IV)                                         & Instrumental Variable                             \\
\rowcolor{Gray}
\cite{ebenstein_particulate_2015}       & 2 Cities, Israel            & Hospital Admissions Due To Lung Illnesses                   & PM10 (air pollutant), Sandstorms (IV)                                              & Instrumental Variable                             \\
\cite{forastiere_assessing_2020}     & Milan, Italy                & Non-Accidental Mortality                                    & Setting PM10 Daily Exposure Levels \textgreater To 40 µg/m³ To 40                  & Generalized Propensity Score                      \\
\rowcolor{Gray}
\cite{giaccherini2021particulate}     & Municipalities, Italy       & Respiratory Hospital Admission                              & PM10, Public Transport Strikes                                     & Difference in Differences                         \\
\cite{godzinski2019short}      & 10 Cities, France           & Emergency Admissions for Upper Respiratory System (Age 0-4) & CO,  Public Transport Strikes                                      & Difference in Differences                         \\
\rowcolor{Gray}
\cite{halliday_vog_2019}       & Hawaii, USA                 & ER Admission for Pulmonary Outcomes                         & PM2.5,  SO2 Emissions From Kilauea Volcano and Wind Direction (IV) & Instrumental Variable                             \\
\cite{he2016effect}            & 34 Urban Districts, China   & Monthly Standardized Mortality Rate                         & PM10, Regulation and Traffic Control Status (IV)                   & Instrumental Variable                             \\
\rowcolor{Gray}
\cite{he2020straw}            & China                       & Monthly Number of Deaths for All-Causes                     & PM2.5, Straw Burning (IV)                                          & Instrumental Variable                             \\
\cite{isphording_pandemic_2021}     & Counties, Germany           & Mortality of Covid-19 Positive Male Patients (Age 80+)      & PM10,  Wind direction (IV)                                         & Instrumental Variable                             \\
\rowcolor{Gray}
\cite{jans_economic_2018}           & Sweden                      & Children Health Care Visits for Respiratory Illness         & PM10, Thermal Inversion (IV)                                       & Reduced-Form                             \\
\cite{jia_is_2019}            & South Korea                 & Mortality Rates for Respiratory and Cardiovascular Diseases & Dusty Days Times China's AQI                                                       & Reduced-Form                                      \\
\rowcolor{Gray}
\cite{kim_sensitivity_2013}            & South Korea                 & Hospital Admissions for Respiratory Illnesses               & PM10 (air pollutant), Average PM10 Level By Date (IV)                              & Instrumental Variable                             \\
\cite{knittel_caution_2016}        & California, USA             & Infant Mortality                                            & PM10, Road Traffic Flow and Weather variables (IV)                 & Instrumental Variable                             \\
\rowcolor{Gray}
\cite{moretti_pollution_2011}        & South California, USA       & Hospital Admissions for Respiratory Illnesses               & O3, Vessel Traffic (IV)                                            & Instrumental Variable                             \\
\cite{mullins_effects_2015}       & Santiago Metropole, Chile   & Cumulative Deaths (age \textgreater 64)                     & PM10, Air quality Alerts                                           & Matching + Difference in Differences              \\
\rowcolor{Gray}
\cite{schlenker_airports_2016}     & California, USA             & Acute Respiratory Hospitalization                           & CO,  Planes Taxi Time (IV)                                         & Instrumental Variable                             \\
\cite{schwartz_estimating_2015}a       & Boston, USA                 & Non-Accidental Mortality                                    & PM2.5, Back Trajectories of PM2.5 (IV)                             & Instrumental Variable                             \\
\rowcolor{Gray}
\cite{schwartz_estimating_2017}       & Boston, USA                 & Non-Accidental Mortality                                    & PM2.5, Height Of Planetary Boundary Layer and Wind Speed (IV)      & Instrumental Variable                             \\
\cite{schwartz_national_2018}      & 135 Cities, USA             & Non-Accidental Mortality                                    & PM2.5, Planetary Boundary Layer, Wind Speed, and Air Pressure (IV) & Instrumental Variable                             \\
\rowcolor{Gray}
\cite{sheldon_impact_2017}     & Singapore                   & Acute Upper Respiratory Tract Infections                    & Pollutant Index, Indonesian Fire Radiative Power (IV)                              & Instrumental Variable                             \\
\cite{williams_short-term_2019}       & USA                         & Asthma Rescue Event                                         & PM2.5                                                             & Poisson fixed-effects models                      \\
\rowcolor{Gray}
\cite{zhong2017traffic}         & Beijing, China              & Ambulance Call Rate for Coronary Heart Problem              & NO2, Number 4 Day (IV)                                             & Instrumental Variable                             \\

\bottomrule
\end{tabular}
\begin{tablenotes}
\scriptsize
\item \textit{Notes}: For each study, we report its location, one of the health outcome analyzed, the independent variables (the air pollutant and in the case of an instrumental variable strategy, the instrument) and the study design.
\end{tablenotes}
\end{threeparttable}
\end{adjustwidth}
\end{table}

    \end{landscape}
    \restoregeometry
    \doublespacing
%-----------------------------------------------------------------------
    
        To evaluate potential statistical power issues in this literature, we first proceed exactly as for the standard epidemiology literature. We compute the statistical power, the exaggeration factor and the probability to get an estimate of the wrong sign for all studies based on hypothetical true effect sizes expressed as decreasing fraction of observed estimates. In \Cref{fig:retro_causal}, each gray line represent the statistical power and average type M error curves of an article. The blue lines represent the average power and exaggeration factor of all causal inference studies.
    
        \begin{figure}[h!]
            \caption{Statistical Power and Type M Error of Causal Inference Studies.}
            \label{fig:retro_causal}
        \centering\includegraphics[width=\linewidth]{2.literature_review/graph_retrodesign_fraction.pdf}
            \caption*{\footnotesize \textit{Notes}: For each causal inference paper, we compute its statistical power and the average type M error for decreasing effect sizes expressed as percentage reduction in observed estimates. Each gray line represents a specific causal inference paper. The blue lines are the average of a metric for all causal inference papers.}
        \end{figure}
    
        If the true effect size of each study was equal to 75\% of the estimate, the median statistical power would be about to 60\% and the median Type M error would be 1.3. In the causal inference literature, at least half of studies have enough statistical power so that statistically significant estimates are not inflated. In \Cref{fig:retro_causal}, we can however see that there is a wide heterogeneity in the robustness of studies to statistical power issues---some of them are relatively well powered while others run quickly into Type M error. A large share of studies in the literature would not have designs with enough statistical power to detect effects of half the size of their observed estimates. In that scenario, the median statistical power would be about 40\% and the median type M error would be 1.8. Overall, this comprehensive retrospective analysis of the literature reveals that some studies are under-powered and could run into type M error. It may help explain why there is a large heterogeneity in effect sizes across articles. 
    
        Again, expressing true effect sizes as decreasing fraction of observed estimates is arbitrary. We also carry out another retrospective analysis where we take as true effect sizes the estimates that would be predicted using non-causal inference methods. We do so for the subset of the 9 instrumental variables that also display estimates in the case when the air pollutant concentration is not instrumented. Two reasons are often advanced in the causal literature to explain the discrepancy between instrumented and non-instrumented estimates: (i) instrumental variables help overcoming omitted variable bias and (ii) if the air pollution is measured with classical error, instrumental variables also reduce the resulting attenuation bias. We think that, for some studies, statistical power issues could also partly explain the observed difference between causal and non-causal methods. In \Cref{tab:retro_iv}, we display the statistical power, the average type M error and the probability to make a type S error for instrumental variable studies. For some studies, the statistical power of the instrumental variable strategy could be extremely low. This results in large type M errors, which magnitude partially close the gap between instrumented and non-instrumented estimates. Given this possibility, future research should carry out quantitative bias analysis to explore the trade-off between using an instrumental variable strategy to overcome omitted variable and attenuation biases and running into a type M error due to low statistical power \citep{rosenbaum2010design, dorie2016flexible, vanderweele2017sensitivity, cinelli2020making}.
    

\begin{table}[h!]
\caption{Retrospective Analysis of Instrumental Variable Papers Where Naive OLS Estimates are Assumed to be True Effect Sizes.}
\begin{threeparttable}
\label{tab:retro_iv}
\begin{tabular}{lccc}
\toprule
\textbf{Paper}                    & \textbf{Power (\%)} & \textbf{Type S Error (\%)} & \textbf{Type M Error} \\ \midrule
\rowcolor{Gray}
\cite{giaccherini2021particulate} & 5                      & 43.3                 & 40.7                          \\
\cite{halliday_vog_2019}   & 6                & 16.6                & 6.9                           \\
\rowcolor{Gray}
\cite{schlenker_airports_2016}   & 7                   & 13.8                 & 6.1                           \\
\cite{moretti_pollution_2011}    & 11                  & 3.7                  & 3.5                           \\
\rowcolor{Gray}
\cite{arceo_does_2016} & 12                   & 2.4                  & 3.1                           \\
\cite{barwick_morbidity_2018}     & 23                   & 0.3                  & 2.1                           \\
\rowcolor{Gray}
\cite{deryugina_mortality_2019}   & 34                  & 0.1                  & 1.7                           \\
\cite{ebenstein_particulate_2015}   & 52                   & 0                    & 1.4                           \\
\rowcolor{Gray}
\cite{schwartz_national_2018}   & 64                     & 0                    & 1.3                           \\ \bottomrule
\end{tabular}
\begin{tablenotes}
\footnotesize
\item \textit{Notes}: For each study based on an instrumental variable strategy, we computed the statistical power, the average type M error and the probability to make a type S error using the non-instrumented estimate as a guess for the true effect size.
\end{tablenotes}
\end{threeparttable}
\end{table}


%------------------------------------------------------------------------------

% SECTION - PROSPECTIVE ANALYSIS OF CAUSAL INFERENCE STRATEGIES 

%------------------------------------------------------------------------------

\section{Prospective Analysis of Causal Inference Methods}

        The review of the standard epidemiology and causal literatures shows that some articles could have produced inflated estimates on the short-term health effects of air pollution. This analysis however does not allow us to clearly identify which parameters of a study influence its statistical power. We therefore implement a prospective analysis to overcome this limitation \citep{gelman_beyond_2014, altoe_enhancing_2020}. We run simulations based on real-data to emulate the main empirical strategies found in the literature. Using real data frees us from the difficult task to model the long-term and seasonal variations in health outcomes but also the specific effects of weather variables such as temperature. In this section, we describe how we implement these simulations. We start by presenting the research designs we wish to simulate, then briefly describe the data we rely on and finally detail how we actually simulate the research designs.
        

    \subsection{Research Designs Simulated}
        
        Several empirical strategies have been implemented to estimate the short-term health effects of air pollution. In our simulations, we try to simulate the main ones found in the literature. We assume below that we are working with data consisting in daily time series of various health outcomes, air pollutant concentrations and weather parameters for a sample of cities.
        
        \paragraph{Standard regression approach.} The standard strategy consists in directly estimating the dose-response between an air pollutant and an health outcome. In the epidemiology literature, researchers often rely on Poisson generalize additive models where the daily count of an health outcome is regressed on the concentration of an air pollutant, while flexibly adjusting for weather parameters, seasonal and long-term variations. Because most causal methods are estimated with linear regression, our simulations are instead based on ordinary least square estimation to approximate the warhorse model used by epidemiologists. We can summarize this approach with the following model:
        
        \begin{equation*}
        \text{Y}_{c,t} = \alpha + \beta\text{P}_{c,t} + \textbf{W}_{c,t}\phi + \textbf{C}_{c,t}\gamma + \epsilon_{c,t}
        \end{equation*}
        
        \noindent where $c$ is the city index and $t$ the daily time index. $\text{Y}_{c,t}$ is the daily count of cases of an health outcome and $\text{P}_{c,t}$ the average daily concentration of an air pollutant. The coefficient $\beta$ measures the short-term effect of an increase in the air pollutant concentration on the health outcome. To address confounding issues, the model adjusts for a set of weather covariates,  $\textbf{W}_{c,t}$, and calendar indicators $\textbf{C}_{c,t}$. The error term is denoted by $\epsilon_{c,t}$.
        

        \paragraph{Instrumental variable approach.} The standard strategy could be prone to omitted variable bias and measurement error. A growing number of articles therefore exploit exogenous variations in air pollution. Most causal inference papers rely on instrumental variable designs where the concentration of an air pollutant is instrumented by thermal inversions \citep{arceo_does_2016}, wind patterns \citep{schwartz_national_2018, deryugina_mortality_2019, isphording_pandemic_2021}, extreme natural events such as sandstorms or volcano eruptions \citep{ebenstein_particulate_2015, halliday_vog_2019}, or variations in transport traffic \citep{moretti_pollution_2011, knittel_caution_2016, schlenker_airports_2016}. This approach can be summarized with a two-stage model where the first stage is:
        
        \begin{equation*}
        \text{P}_{c,t} = \delta + \theta\text{Z}_{c,t} + \textbf{W}_{c,t}\phi + \textbf{C}_{c,t}\gamma + e_{c,t}
        \end{equation*}
        
        \noindent where $\text{Z}_{c,t}$ is the instrumental variable. The second stage is then:
        
        \begin{equation*}
        \text{Y}_{c,t} = \alpha + \beta\widehat{\text{P}}_{c,t} + \textbf{W}_{c,t}\psi + \textbf{C}_{c,t}\lambda + \epsilon_{c,t}
        \end{equation*}
        
        \noindent where $\widehat{\text{P}}_{c,t}$ is the exogenous variation in an air pollutant predicted by the instrument. The causal effect measured by this approach is a weighted average of per-unit causal responses to an increase in the concentration of an air pollutant \citep{angrist1995two}.

        \paragraph{Reduced-form approach.} A subset of articles however restrain from instrumenting the concentrations of an air pollutant with exogenous shocks, but instead choose to directly estimate the relationship between the health outcome and the shocks. This why we call this approach a reduced-form analysis. The articles concerned by this approach focus on public transport strikes \citep{bauernschuster_when_2017, godzinski2019short, giaccherini2021particulate}. Their empirical strategy can be summarized with the following model:
        
        \begin{equation*}
        \text{Y}_{c,t} = \alpha + \beta\text{Z}_{c,t} + \textbf{W}_{c,t}\phi + \textbf{C}_{c,t}\gamma + \epsilon_{c,t}
        \end{equation*}
        
        \noindent The coefficient $\beta$ captures a type of intention-to-treat effect.
        
        \paragraph{Regression-discontinuity design approach.} The final empirical strategy found in the literature consists in measuring the effects of air quality alerts with a regression-discontinuity design \citep{chen_effect_2018}. The approach is summarized with the model:
        
        \begin{equation*}
        g\{E(\text{Y}_{c,t})\} = \beta_{0} + \beta_{1}(I_{c,t}-l) + \beta_{2}E_{c,t} + \beta_{3}(I_{c,t}-l)\times E_{c,t} + \epsilon_{c,t}
        \end{equation*}
        
        \noindent where $g(.)$ is a generic link function, I$_{c,t}$ is the daily air quality index of a city, $E_{c,t}$ is the indicator for the eligibility to issue an air quality alert whose threshold is $l$. This approach estimates the intention-to-treat effect of air quality alerts. It is important to note that it can capture both the effect due to a subsequent decrease in air pollution due to traffic restriction policies and the effect caused by inhabitants' avoidance behavior.
        
        \subsection{Data}
        
        Our simulation exercises are based on a subset of the US National Morbidity, Mortality, and Air Pollution Study (NMMAPS). The dataset has been exploited in several major studies of the early 2000s to measure the short-term effects of ambient air pollutants on mortality outcomes \citep{peng2008statistical}. It is publicly available and allows us to work with increasing sample sizes for our simulations. Specifically, we extracted daily data on 68 cities over the 1987-1997 period, which represent 4,018 observations per city, for a total sample size of 273,224 observations. For each city, the average temperature (C°), the standardized concentration of carbon monoxide (CO), and mortality counts for several causes are recorded. We choose to work with CO as it is the air pollutant measured in most cities over the period. Less than 5\% of carbon monoxide concentrations and average temperature readings are missing in the initial data set and we impute them using the chained random forest algorithm provided by the \texttt{missRanger} package \citep{mayer_missranger_2019}.
        
        \subsection{Simulations Set-Up}
        
        \paragraph{Simplifying assumptions.} To only capture the specific issues arising due to low statistical power, we make several simplifications to make sure that all the assumption of empirical strategies are met. First, in all research designs, there is no bias due to unmeasured confounders or measurement errors. All models retrieve on average the true value of the treatment effect we set in the data. Second, for instrumental variable and reduced-form strategies, we only simulate exogenous shocks that are binary (e.g. the occurrence of a thermal inversion or not). They are randomly allocated. Third, for the regression discontinuity approach, we only model sharp designs where an air quality is always activated above a randomly chosen threshold.
     
        \paragraph{Two approaches for simulating research designs.} We take two different approaches to simulate the research designs. 
        
        For the reduced-form and regression discontinuity designs, we follow the Neyman-Rubin causal framework by creating a Science table \citep{rubin_estimating_1974}. The recorded value of a health outcome in the dataset represent the potential outcome Y$_{c,t}$(0) of that day $t$ in city $c$ when it is not exposed to the treatment denoted by W$_{c,t}$. It is equal to 1 when a treatment occurs and 0 otherwise. To create the counterfactuals Y$_{c,t}$(1), we add a treatment effect drawn from a Poisson distribution whose parameter correspond to the magnitude of the treatment. We then randomly draw the treatment indicators W$_{t,c}$ for exogenous shocks or air quality alerts. For reduced-form strategies, the treatment status of each day is drawn from a Bernoulli distribution with parameter equal to the proportion of exogenous shocks desired. For air pollution alerts, we randomly draw a threshold from a uniform distribution and select a bandwidth such that it yields the correct proportion of treated observations. We finally express the observed values Y$^{obs}$ of potential outcomes according to the treatment assignment: Y$^{obs}_{c,t}$ = (1-W$_{c,t}$)$\times$Y$_{c,t}$(0) + W$_{c,t}$$\times$Y$_{c,t}$(1).
      
        For the standard regression and the instrument variable strategies, we rely on a model-based approach. For the standard regression strategy, we first estimate the following statistical model on our data:
        
        \begin{equation*}
        \text{Y}_{c,t} = \alpha + \beta\text{Z}_{c,t} + \textbf{W}_{c,t}\phi + \textbf{C}_{c,t}\gamma + \epsilon_{c,t}
        \end{equation*}
        
        \noindent We then predict new observations of a $\text{Y}_{c,t}$ using the estimated coefficients of the model ($\hat{\beta}$, $\hat{\phi}$, and $\hat{\gamma}$) and by adding noise drawn from the residuals distribution $\widehat{\epsilon_{c,t}}$ \citep{peng2006model}. We modify the slope of the dose-response relationship by changing the value of the air pollution coefficient $\beta$. For the instrumental variable strategy, we use the same method as for the standard regression approach but first modify observed air pollutant concentrations $\text{P}_{c,t}$ according to the desired effect size $\theta$ of the randomly allocated instrument: 
        
        \begin{equation*}
        \text{P}_{c,t} = \text{P}_{c,t} + \theta\text{Z}_{c,t}
        \end{equation*}
        
        \noindent The allocation of each day to an exogenous shock is drawn from a Bernoulli distribution with parameter equal to the proportion of exogenous shocks desired. We then estimate a two-stage least squares model, modify the coefficient for the effect of the air pollutant on an health outcome, and finally generate the fake observations of the health outcome using coefficients of the two-stage least squares model and noise drawn from the residuals of the estimated model.

        \paragraph{General procedure.} Our simulation procedure therefore follows 7 main steps:

        \begin{senumerate}
            \item Randomly draw a study period and a sample of cities.
            \item For instrumental variable, reduced-form and regression-discontinuity designs, randomly allocate days to binary exogenous shocks/air quality alerts.
            \item Add the treatment effect size of interest.
            \item Run the model of the empirical strategy.
            \item Store the point estimate of interest and its standard error.
            \item Repeat the procedure 1000 times.
            \item Finally compute the statistical power, the exaggeration ratio of statistically significant estimates and the probability that they are of the wrong sign.
        \end{senumerate}
        
        \paragraph{Varying parameters.} To understand which parameters affect statistical power issues, we can change one aspect the research design while keeping other parameters values constant. We consider in our simulations four main parameters. First, by drawing a different number of cities and changing the length of the period, we vary the sample size. Second, we change the effect size of air pollution or an exogenous shock on an health outcome. Third, we also allocate increasing proportions of exogenous shocks/air quality alters. Fourth, we can run similar simulations but for different health outcomes with small or large number of cases per day.
        
        \subsection{Simulating Flagship Studies.} Our simulations based on real data help explore the consequences of varying each parameter on statistical power issues. We could however make the results of our simulations even more informative by setting realistic values for the four parameters of a research design. We therefore also try to reproduce three flagship studies using our own dataset.
        

%------------------------------------------------------------------------------

% SECTION - RESULTS

%------------------------------------------------------------------------------

\section{Results}

In this section, we first explore how statistical power is related to the treatment effect size, the number of observations, the proportion of treated units and the distribution of the health outcome. We then try to replicate the design of flagship publications to highlight their potential weaknesses with respect to low statistical power issues. 

    \subsection{Evolution of Power, Type M and S Errors with Study Parameters}

            First, we analyze how statistical power, type M and S errors are affected by the value of different study parameters. To do so, we set baseline values for these parameters and vary the value of each of them one by one. This enables us to get a sense of the impact of each parameter, other things being equal. The baseline parameters are such that:
            
            \begin{itemize}
                \item The sample size is equal to 100,000 observations (2500 days $\times$ 40 cities).
                \item The effect size of air pollution or an exogenous shock is equal to a 1\% relative increase in an health outcome.
                \item The proportion of exogenous shocks represents 50\% of observations. For air pollution alerts analyzed with regression discontinuity designs, we choose a smaller proportion of treated units: 10\%.
                \item The health outcome is the total daily number of non-accidental deaths. It is the health outcome with the largest average number of counts---the average daily mean is equal to 23 cases.
            \end{itemize}
            
            \noindent For all statistical models, we adjust for temperature, temperature squared, city  and calendar (weekday, month, year, month$\times$year) fixed effects. We also repeat the simulations for a smaller sample size of 10,000 observations.
        
        \subsubsection*{Sample Size}

            As shown in \Cref{fig:evol_typeM_n_days}, we unsurprisingly find that, for all identification methods, statistical power increases and type M error decreases with the number of observations.
            
            \begin{figure}[h!]
                \caption{Evolution of Type M Error against Sample Size.}
                \label{fig:evol_typeM_n_days}
                \centering\includegraphics[width=\linewidth]{3.simulations/graph_type_m_sample_size.pdf}
                \caption*{\footnotesize \textit{Notes}: The true effect size is a 1\% relative increase in the health outcome. The health outcome used in the simulations is the total number of non-accidental deaths. The proportion of exogenous units is 50\% for instrumental variable and reduced-form designs.}
            \end{figure}
            
            \noindent Yet, statistical power and type M error issues arise even for a large number of observations. For a sample size of 40,000 observations, an instrumental variable strategy would only have a statistical power of 54\% and would overestimate the true effect by a factor of 1.4. On the contrary, a standard regression strategy is much less prone to power issues than the instrumental variable strategy. This is explained by the fact that the variance of the two stage least-square estimator is larger than the variance of the ordinary least square estimator. In our simulations, we also note that, for all identification method, Type S error is not a problem for any sample sizes. 
            
        \subsubsection*{Effect Size}
        
            The second unsurprising result of our simulations is that the larger the effect size, the larger the power and the lower type M and S errors are. With our advantageous baseline parameters, statistical power issues however start to appear in instrumental variable and regression discontinuity designs for effect sizes below 1\%. For instance, for an effect of 0.5\%, the average type M error is about 1.7. Such effect sizes are similar to those sometimes found in the standard epidemiology literature. As for results on sample sizes, standard regression and reduced-form strategies suffer less from power issues, even for small effects.

        \subsubsection*{Proportion of Exogenous Shocks}
        
            The link between the proportion of exogenous shocks and statistical power might be less known to researchers. In \Cref{fig:evol_power_p_treat}, we see that the statistical power increases with larger proportions of treated units for instrumental variable, regression discontinuity and reduced-form designs. As in the case of randomized controlled trials, the precision of studies will be maximized when half of the observations are exposed to the treatment of interest.
            
            \begin{figure}[h!]
                \caption{Evolution of Statistical Power with the Proportion of Exogenous Shocks.}
                \label{fig:evol_power_p_treat}
                \centering\includegraphics[width=\linewidth]{3.simulations/graph_power_prop_exogenous_shocks.pdf}
                \caption*{\footnotesize \textit{Notes}: The true effect size is a 1\% relative increase in the health outcome. The health outcome used in the simulations is the total number of non-accidental deaths.}
            \end{figure}
            
           \noindent Conversely, as shown in \Cref{fig:evol_typeM_p_treat}, the average Type M error increases as the proportion of exogenous shocks decreases.
            
            \begin{figure}[h!]
                \caption{Evolution of Type M Error with the Proportion of Exogenous Shocks.}
                \label{fig:evol_typeM_p_treat}
                \centering\includegraphics[width=\linewidth]{3.simulations/graph_type_m_prop_exogenous_shocks.pdf}
                \caption*{\footnotesize \textit{Notes}: The true effect size is a 1\% relative increase in the health outcome. The health outcome used in the simulations is the total number of non-accidental deaths.}
            \end{figure}
            
            \noindent Air pollution alerts, thermal inversion or transportation strikes are however rare events. They can represent less than 5\% of the observations in some studies. With a dataset of 10,000 observations, the average type M error is 2.7 for reduced-form strategies. The causal inference literature might therefore be particularly prone to type M error due to a very low proportion of treated units, even though sample sizes are often large.

        \subsubsection*{Average Count of Cases of the Health Outcome}
            
            Perhaps less known to economists than the influence of sample and effect sizes, the average count of cases also critically affects statistical power. For instance, a 1\% increase in the number of deaths in a setting where there are only 2 deaths per day corresponds to rare additional deaths that might therefore be more difficult to detect. To simulate situations with various number of cases, we consider three different outcome variables, with different counts of cases: the total number of non-accidental deaths (daily mean $\simeq$ 23), the total number of respiratory deaths (daily mean $\simeq$ 2) and the number of chronic obstructive pulmonary disease cases for people aged between 65 and 75 (daily mean $\simeq$ 0.3). Using baseline parameters and in the case of the large dataset, we find that statistical power is close to 100\% when empirical strategies target a 1\% increase in the total number of non-accidental deaths. However, statistical power quickly drops when the average count of cases decreases. For instance, an instrumental variable strategy has only 16\% of statistical power to detect an increase by 1\% in respiratory deaths. The average type M error is then equal to 2.4. For chronic obstructive pulmonary deaths, the situation is even worse, with an average type M error of 5.9. Studies with a small count of cases may therefore lead to extreme statistical power issues.
            
        \subsubsection*{Issues Specific to the Instrumental Variable Design}
            
            For instrumental variable strategies, we also analyze how the statistical power is affected by the strength of the instrument. In our simulations, we define the strength of the instrument as the standardized effect size on the air pollutant concentration. A strength equals to 0.2 means that the instrument increases the concentration by 0.2 standard deviation. 
            
                \begin{figure}[h!]
                \caption{Evolution of Type M Error with the Strength of the Instrumental Variable.}
                \label{fig:iv_issues}
                \centering\includegraphics[width=\linewidth]{3.simulations/graph_power_issues_iv.pdf}
                \caption*{\footnotesize \textit{Notes}: The true effect size is a 1\% relative increase in the health outcome. The health outcome used in the simulations is the total number of non-accidental deaths. Half of the observations are exposed to exogenous shocks. The strength of the instrumental variable is its effect in standard deviation on the air pollutant concentration.}
            \end{figure}
            
            As shown in \Cref{fig:iv_issues}, we find that statistical power collapses and type M error soars when the instrument's strength decreases. Importantly, this issue arises for rather large instrument's strengths. Even in the case of the large data set with 100,000 observations, an instrumental variable's strength of 0.2, and effect size of a 1\% increase in the health outcome, statistical power is only 23\% and the average type M error is 2. This statistical power issue arises for a \textit{F}-statistics of 1278!  A large \textit{F}-statistic could therefore hide a weak instrumental variable that results in a low statistical power.
            
        \subsection{Simulating Flagship Studies}

            The simulation results of the previous section help build the intuition for the parameters influencing the statistical power of studies. Yet, they represent an ideal setting, with relatively large sample size, proportion of treated units, outcome counts and instrumental variable strength. These parameters may not perfectly represent actual studies. For each causal inference method, we therefore consider a realistic set of parameters based on examples from the literature. We then vary the value of key parameters one by one in order to see what could be changed in each study to avoid running into power issues.
        
        \subsubsection*{Public Transport strikes}
        
            Public transport strikes are unique but rare exogenous events where air pollution increases. Even in a large data set, with several cities and a long study period, the proportion of treated days might be very small. For instance, \cite{bauernschuster_when_2017} investigate the effect of public transport strikes on air pollution and emergency admission in the five biggest German cities over a period of 6 years. The sample size of the study is equal to 11,000 observations but there are only 45 1-day strikes. This study could be prone to statistical power issues since the proportion of treated units is 0.4\%. We thus try to simulate with our data a similar design. In our baseline simulation, we set as the true effect size the point estimate found by \cite{bauernschuster_when_2017}: days with strikes see an 11\% relative increase in the health outcome of interest. The average count of cases for our health outcome---the total number of respiratory deaths--is however 3 times larger than the one in their study, which is equal to 0.69.
            
            In the baseline scenario, we find that the statistical power is only 15\% and the average type M error is 2.7. If the researchers had looked at the effect for an health outcome with an average of 23 cases per day, there would however be no statistical power issues. The effect size found by the authors could nonetheless be argued to be a very large increase in an health outcome. If the true effect was only 5\% and the average count of the health outcome was 23, there would still be a substantial risk to overestimate statistically significant estimates by a factor of 1.8. Estimating the effects of rare exogenous events on health outcomes with few cases could be therefore difficult.
        
        \subsubsection*{Air pollution Alerts}
        
            Air pollution alerts are also rare events. Contrary to public transport strikes or thermal inversions, their effects are estimated using regression discontinuity design. Only observations closed to the air quality threshold are included in the analysis. As a consequence, the effective sample size may end up being particularly small. For instance, in \cite{chen_effect_2018}, while the initial sample size is equal to 3652 observations, the effective sample size is only of 143 (100 control observations and 43 treated ones). The proportion of treated observation is 1.2\%. With our data, we try to approximate the setting of \cite{chen_effect_2018}. In the baseline scenario, we sample one city with a time period of 3652 days and randomly allocate the treatment to 1.2\% of observations. We also consider a true effect size of 12\%, as found in the study. The average number of cases of their health outcome is 26 cases per day. In our simulations, we use the total number of non-accidental deaths as our outcome variable since the daily mean is equal to 23 deaths.
            
           In the baseline scenario, we find that the statistical power is only 10\% and the average type M error is 4.6. If we consider smaller true effect sizes, type M error shoots up and power collapses. As a consequence, care must be taken when interpreting estimates from studies with such a small sample size and few air quality alerts. Of course, one could always argue that air quality alerts have large and protective effects on health outcomes if individuals adopt avoidance behavior.
            
        \subsubsection*{Instrumenting Air Pollution}
        
            Finally, we investigate the most common strategies used in the causal inference literature, which are based on instrumental variables. These papers often rely on very large datasets. For instance \cite{schwartz_national_2018} gathered 591,570 observations (135 cities with a length of study of approximately 4382 days). In this study, air pollution is instrumented with a complex mix of variables and we cannot easily observe the proportion of treated units. The effect size found by the authors is equal to a 1.5\% relative increase in an health outcome with an average daily number of cases equal to 23. In our simulations, we therefore assume that half of the observations are exposed to exogenous shocks. We only vary the strength of the instrument and use the total number of non-accidental deaths as the outcome variable. Our data set being smaller than the one used in the study, we only consider 2500 days and 40 cities.
            
            If the instrumental variable increases air pollution concentration by 0.5 standard deviation, we find a statistical power of nearly 100\% and an average type M error of 1. Yet, for smaller values of the instrument's strength, statistical power rapidly decreases. For an instrument's strength of 0.2, the statistical power is 48\% and the average type M error is 1.4. For a strength of 0.1, power is only 16\% and the average type M error is 2.6. In these two scenarios, the values of the \textit{F}-statistic remain extremely large, with respective values equal to 1287 and 320. A large \textit{F}-statistic could be a poor indicator of statistical power issues.

%------------------------------------------------------------------------------

% SECTION - DISCUSSION

%------------------------------------------------------------------------------

\section{Discussion}

\epigraph{``\textit{I think that when we know that we actually do live in uncertainty, then we ought to admit it.}"}{--- Richard P. Feynman}

        Our findings should make us worried about statistical power issues when we are trying to estimate the acute health effects of air pollution. Our retrospective analysis of the literature suggests that under-powered studies with inflated effect sizes could be a real issue both in the standard epidemiology and the causal inference literatures. We thus recommend to adopt retrospective calculations since they are very easy to implement and force us to reflect on the range of plausible effect sizes we are trying to estimate. 

        Unfortunately, a retrospective analysis will not help us understand which parameters of the research design  that influence the statistical power of our studies. Our prospective analysis, using simulations based on real-data, fills this gap and leads to issue four warnings. First, sample size matters for all causal inference methods but especially for the regression-discontinuity design applied to air pollution alerts. Given the sample size its entails, we advise researchers to interpret findings with extra care as the inflation of statistically significant estimates can be extremely large, even when we assume that true effect sizes are large. Second, despite their large sample sizes, when we exploit rare exogenous shocks such as transport strikes, we should be aware that the small proportion of exogenous shocks observed in our studies can lead to a dramatically low statistical power. Third, although it is well-known that two-stage least square estimates are inherently less precise than ordinary least square estimates, it also makes instrumental variable strategies more prone to power issues. If one thinks that omitted variable and attenuation biases are small, the benefits of using an instrumental variable strategy could be questioned. The trade-off between targeting an unbiased estimate with causal inference methods and the risk of running into a type M error could be a fruitful area of research for quantitative bias analysis \citep{rosenbaum2010design, dorie2016flexible, vanderweele2017sensitivity, cinelli2020making}. Fourth, the power of all research designs in the literature is driven by the average count of the health outcome. Many articles investigate the acute effects of air pollution for specific groups such as children and the elderly. In such settings, there is potentially a huge risk to make a type M error, even with large sample sizes. While they are more involved than a retrospective analysis, simulating the research design we want to implement is the best way to assess if it could suffer from statistical power issues. Our simulation codes in the replication material provide a template to run such prospective analysis.

        On top of these specific considerations, we think that the literature would benefit from adopting a different view towards statistically insignificant results \citep{ziliak_cult_2008, wasserstein2016asa, mcshane2019abandon}. The null hypothesis testing framework remains very strong in the field, especially for causal inference papers, since nearly all of them dichotomize evidence using the 5\% significance threshold \citep{greenland2017invited}. This statistical significance filter leads to publication bias and is at the very heart of the inflation of statistically significant estimates in under-powered studies \citep{amrhein2019statistical, gelman2020regression, romer2020praise}. Even if we could not improve the statistical power of our studies, the distribution of the acute health effects of air pollution could be more accurate if statistically insignificant results were not kept in the file drawer \citep{hernan2022causal}.
        
        Finally, our results show that a credible identification strategy does not necessarily lead to a correct estimation of the actual true effect \citep{young2019consistency}. When we qualify estimates as "statistically significant", there is often much more uncertainty lying behind, an uncertainty that should be computed and embraced to better help policy-makers evaluate the adverse effects of air pollution. We are convinced that prospective and retrospective power analyses can help us design better studies and improve the interpretation of their results.


\setlength\bibsep{0pt}
\bibliographystyle{plainnat}
\bibliography{air_pollution.bib}

\end{document}
